\documentclass[aspectratio=169,11pt]{beamer}

\usepackage[T1]{fontenc}
\usepackage[utf8]{inputenc}
\usepackage{lmodern}
\usepackage{microtype}
\usepackage{booktabs}
\usepackage{tabularx}
\usepackage{array}
\usepackage{ragged2e}
\usepackage{xcolor}
\usepackage{hyperref}

% -----------------------------------------------------------------------------
% Identitas dan palet StudentPro
% -----------------------------------------------------------------------------
\definecolor{StudentNavy}{HTML}{0B1F3A}
\definecolor{StudentBlue}{HTML}{1769E0}
\definecolor{StudentCyan}{HTML}{18B6C9}
\definecolor{StudentGold}{HTML}{F4B942}
\definecolor{StudentInk}{HTML}{182230}
\definecolor{StudentMist}{HTML}{EFF5FB}
\definecolor{StudentSoft}{HTML}{D9E7F5}
\definecolor{StudentGray}{HTML}{607086}
\definecolor{StudentWhite}{HTML}{FFFFFF}

\hypersetup{
  colorlinks=true,
  urlcolor=StudentBlue,
  linkcolor=StudentBlue
}

\setbeamercolor{normal text}{fg=StudentInk,bg=StudentWhite}
\setbeamercolor{structure}{fg=StudentBlue}
\setbeamercolor{frametitle}{fg=StudentNavy,bg=StudentWhite}
\setbeamercolor{title}{fg=StudentWhite}
\setbeamercolor{subtitle}{fg=StudentSoft}
\setbeamercolor{author}{fg=StudentWhite}
\setbeamercolor{date}{fg=StudentSoft}
\setbeamercolor{block title}{fg=StudentWhite,bg=StudentBlue}
\setbeamercolor{block body}{fg=StudentInk,bg=StudentMist}
\setbeamercolor{alerted text}{fg=StudentBlue}

\setbeamerfont{title}{series=\bfseries,size=\fontsize{30}{35}\selectfont}
\setbeamerfont{subtitle}{size=\fontsize{16}{20}\selectfont}
\setbeamerfont{author}{series=\bfseries,size=\fontsize{16}{19}\selectfont}
\setbeamerfont{date}{size=\fontsize{12}{15}\selectfont}
\setbeamerfont{frametitle}{series=\bfseries,size=\fontsize{23}{27}\selectfont}
\setbeamerfont{normal text}{size=\fontsize{16}{20}\selectfont}
\setbeamerfont{block title}{series=\bfseries,size=\fontsize{15}{18}\selectfont}
\setbeamerfont{block body}{size=\fontsize{14}{18}\selectfont}
\setbeamerfont{footline}{size=\fontsize{8}{10}\selectfont}

\setbeamertemplate{navigation symbols}{}
\setbeamertemplate{items}[circle]
\setbeamertemplate{itemize item}{\color{StudentCyan}\raisebox{0.12ex}{\scriptsize$\blacktriangleright$}}
\setbeamertemplate{itemize subitem}{\color{StudentBlue}\scriptsize$\bullet$}
\setlength{\leftmargini}{1.2em}
\setlength{\leftmarginii}{1.1em}

\setbeamertemplate{frametitle}{%
  \nointerlineskip
  \begin{beamercolorbox}[wd=\paperwidth,leftskip=0.7cm,rightskip=0.7cm,sep=0.18cm]{frametitle}
    \usebeamerfont{frametitle}\insertframetitle\par
    \color{StudentCyan}\rule{\dimexpr\paperwidth-1.8cm\relax}{0.8pt}
  \end{beamercolorbox}
}

\setbeamertemplate{footline}{%
  \leavevmode
  \hbox{%
    \begin{beamercolorbox}[wd=.64\paperwidth,ht=2.7ex,dp=1.2ex,leftskip=.7cm]{author in head/foot}%
      \color{StudentGray}\usebeamerfont{footline}StudentPro \textbar\ SMAN 1 Muntilan
    \end{beamercolorbox}%
    \begin{beamercolorbox}[wd=.34\paperwidth,ht=2.7ex,dp=1.2ex,center]{date in head/foot}%
      \color{StudentBlue}\usebeamerfont{footline}Bab 1 \textbar\ \insertframenumber/\inserttotalframenumber
    \end{beamercolorbox}%
  }
}

\newcommand{\kicker}[1]{%
  {\color{StudentCyan}\bfseries\fontsize{12}{14}\selectfont\MakeUppercase{#1}}\par\vspace{0.15cm}
}

\newcommand{\claim}[1]{%
  {\color{StudentNavy}\bfseries\fontsize{22}{27}\selectfont #1}\par
}

\newcommand{\source}[1]{%
  \vfill
  {\color{StudentGray}\fontsize{7.5}{9}\selectfont Sumber: #1}\par
}

\newcommand{\metric}[2]{%
  {\color{StudentBlue}\bfseries\fontsize{27}{30}\selectfont #1}\par
  \vspace{0.05cm}
  {\color{StudentGray}\fontsize{12}{15}\selectfont #2}
}

\title{Pendahuluan dan Tujuan\\Pembimbingan UTBK 2027}
\subtitle{Bab 1 \textbar\ Program SMAN 1 Muntilan berbasis data dan LMS StudentPro}
\author{Balya Rochmadi}
\date{Tahun Ajaran 2026/2027}

\begin{document}

% Slide 1 ---------------------------------------------------------------------
\begingroup
\setbeamercolor{background canvas}{bg=StudentNavy}
\begin{frame}[plain]
  \vspace{0.75cm}
  \kicker{Rencana Pembimbingan}
  \vspace{0.1cm}
  {\usebeamerfont{title}\color{StudentWhite}\inserttitle\par}
  \vspace{0.45cm}
  {\usebeamerfont{subtitle}\color{StudentSoft}\insertsubtitle\par}
  \vfill
  {\color{StudentGold}\rule{2.2cm}{2.5pt}}\par
  \vspace{0.22cm}
  {\usebeamerfont{author}\color{StudentWhite}\insertauthor}\par
  {\usebeamerfont{date}\color{StudentSoft}\insertdate}
  \vspace{0.45cm}
\end{frame}
\endgroup

% Slide 2 ---------------------------------------------------------------------
\begin{frame}{UTBK menuntut lebih dari hafalan}
  \begin{columns}[T,onlytextwidth]
    \begin{column}{0.57\textwidth}
      \kicker{Titik Berangkat}
      \claim{Keberhasilan ditentukan oleh penalaran, literasi, ketepatan strategi, dan kendali waktu.}
      \vspace{0.35cm}
      {\fontsize{15}{19}\selectfont
      Framework resmi 2026 menempatkan peserta dalam dua kelompok tes, tujuh subtes, dan waktu total 195 menit tanpa istirahat.}
    \end{column}
    \begin{column}{0.37\textwidth}
      \vspace{0.18cm}
      \begin{block}{Implikasi untuk 2027}
        Pembimbingan harus melatih proses berpikir dan keputusan saat tertekan waktu, bukan sekadar menambah jumlah soal.
      \end{block}
    \end{column}
  \end{columns}
\end{frame}

% Slide 3 ---------------------------------------------------------------------
\begin{frame}{Tiga dasar penyusunan program 2027}
  \kicker{Keputusan Berbasis Bukti}
  \vspace{0.05cm}
  {\fontsize{13}{16}\selectfont
  \renewcommand{\arraystretch}{1.08}
  \begin{tabularx}{\textwidth}{>{\bfseries\color{StudentBlue}}p{0.23\textwidth} >{\RaggedRight\arraybackslash}X}
    \toprule
    Kerangka resmi & Acuan kompetensi, jenis tes, dan tuntutan pengerjaan. \\
    Evaluasi 2026 & Dasar menemukan pola yang perlu dipertahankan atau diperbaiki. \\
    Data LMS & Dasar menentukan tindak lanjut belajar setiap siswa. \\
    \bottomrule
  \end{tabularx}}
  \vspace{0.22cm}
  {\color{StudentNavy}\bfseries\fontsize{13}{16}\selectfont Fakta resmi, data internal, dan proyeksi diberi label yang berbeda.}
\end{frame}

% Slide 4 ---------------------------------------------------------------------
\begin{frame}{LBI berubah: skoring dan muatan}
  \kicker{Arah Penyesuaian 2027}
  {\fontsize{13}{16}\selectfont
  \renewcommand{\arraystretch}{1.10}
  \begin{tabularx}{\textwidth}{>{\bfseries\color{StudentBlue}}p{0.23\textwidth} >{\RaggedRight\arraybackslash}X}
    Hasil 2026 & Skor LBI dirinci menjadi konteks \textbf{saintek} dan \textbf{soshum}. \\
    Rencana 2027 & Muatan \textbf{Tes Wawasan Kebangsaan} direncanakan masuk ke LBI; bentuk final menunggu publikasi resmi. \\
    Kesiapan LMS & Soal dan skor LBI ditandai per konteks, termasuk wawasan kebangsaan. \\
  \end{tabularx}}
\end{frame}

% Slide 5 ---------------------------------------------------------------------
\begin{frame}{Kebutuhan siswa tidak sama}
  \begin{columns}[T,onlytextwidth]
    \begin{column}{0.46\textwidth}
      \kicker{Kondisi Umum}
      \claim{Materi yang sama tidak selalu menghasilkan kemajuan yang sama.}
    \end{column}
    \begin{column}{0.49\textwidth}
      \begin{itemize}
        \item Ada siswa yang lemah pada konsep dasar.
        \item Ada yang memahami materi, tetapi lambat mengambil keputusan.
        \item Ada yang stabil di latihan, tetapi menurun saat tryout penuh.
        \item Ada yang membutuhkan strategi pilihan program studi yang lebih realistis.
      \end{itemize}
    \end{column}
  \end{columns}
  \vspace{0.35cm}
  \begin{beamercolorbox}[rounded=true,sep=0.23cm,wd=\textwidth]{block body}
    \centering\bfseries Program 2027 harus mampu mengubah hasil tes menjadi intervensi yang spesifik.
  \end{beamercolorbox}
\end{frame}

% Slide 6 ---------------------------------------------------------------------
\begin{frame}{Tujuan pembimbingan UTBK 2027}
  \kicker{Hasil yang Diharapkan}
  \begin{enumerate}
    \item Memetakan kemampuan awal siswa sampai tingkat subtes dan indikator.
    \item Menyediakan materi yang terstruktur, bertahap, dan mudah diakses.
    \item Membentuk kebiasaan mengerjakan soal secara akurat dalam batas waktu.
    \item Memberikan remedial dan pengayaan berdasarkan bukti hasil belajar.
    \item Mendampingi pemilihan program studi dengan pertimbangan kemampuan dan peluang.
    \item Menjaga konsistensi belajar melalui monitoring sekolah dan StudentPro.
  \end{enumerate}
\end{frame}

% Slide 7 ---------------------------------------------------------------------
\begin{frame}{Indikator proses menopang target akhir}
  \kicker{Ukuran Keberhasilan}
  \begin{columns}[T,onlytextwidth]
    \begin{column}{0.28\textwidth}
      \metric{70\%}{target siswa peserta program diterima di perguruan tinggi tujuan melalui jalur yang direncanakan.}
    \end{column}
    \begin{column}{0.66\textwidth}
      \begin{itemize}
        \item Seluruh siswa mengikuti diagnostik awal.
        \item Progres materi dan latihan tercatat secara konsisten.
        \item Nilai tryout menunjukkan tren yang dapat dijelaskan.
        \item Kelemahan indikator memperoleh remedial yang terukur.
        \item Setiap keputusan jurusan didukung data akademik dan konsultasi.
      \end{itemize}
    \end{column}
  \end{columns}
  \vspace{0.25cm}
  {\color{StudentGray}\fontsize{10.5}{13}\selectfont Target 70\% adalah sasaran internal StudentPro, bukan ketentuan SNPMB.}
\end{frame}

% Slide 8 ---------------------------------------------------------------------
\begin{frame}{Lima prinsip pembimbingan ke depan}
  \kicker{Cara Kerja Program}
  {\fontsize{13.5}{17}\selectfont
  \begin{itemize}
    \item \textbf{Diagnosis dahulu.} Materi dipilih setelah kemampuan awal diketahui.
    \item \textbf{Sedikit tetapi konsisten.} Latihan rutin selalu disertai evaluasi.
    \item \textbf{Umpan balik cepat.} Kesalahan dibahas sebelum menjadi pola.
    \item \textbf{Simulasi autentik.} Tryout melatih waktu, stamina, dan keputusan.
    \item \textbf{Tindak lanjut individual.} Data siswa menentukan intervensi berikutnya.
  \end{itemize}}
\end{frame}

% Slide 9 ---------------------------------------------------------------------
\begin{frame}{StudentPro menjadi tulang punggung program}
  \kicker{Satu Siklus Pembelajaran}
  \begin{columns}[T,onlytextwidth]
    \begin{column}{0.23\textwidth}
      \textbf{Belajar}\par
      \vspace{0.12cm}
      {\fontsize{12.5}{15.5}\selectfont Materi per subtes, capaian membaca, dan akses semi-offline.}
    \end{column}
    \begin{column}{0.23\textwidth}
      \textbf{Berlatih}\par
      \vspace{0.12cm}
      {\fontsize{12.5}{15.5}\selectfont Kuis bertingkat, bank soal, paket subtes, dan tryout penuh.}
    \end{column}
    \begin{column}{0.23\textwidth}
      \textbf{Menganalisis}\par
      \vspace{0.12cm}
      {\fontsize{12.5}{15.5}\selectfont Skor, waktu, ketuntasan, serta rincian LBI saintek dan soshum.}
    \end{column}
    \begin{column}{0.23\textwidth}
      \textbf{Menindaklanjuti}\par
      \vspace{0.12cm}
      {\fontsize{12.5}{15.5}\selectfont Remedial, pengayaan, konsultasi, dan rekomendasi belajar berikutnya.}
    \end{column}
  \end{columns}
  \vspace{0.15cm}
  {\color{StudentBlue}\bfseries\fontsize{12}{15}\selectfont Data selalu berakhir pada keputusan: lanjut, remedial, pengayaan, atau konsultasi.}
\end{frame}

% Slide 10 --------------------------------------------------------------------
\begin{frame}{Setiap pihak memegang peran berbeda}
  \kicker{Akuntabilitas Bersama}
  {\fontsize{13}{16}\selectfont
  \renewcommand{\arraystretch}{0.98}
  \begin{tabularx}{\textwidth}{>{\bfseries\color{StudentBlue}}p{0.20\textwidth} >{\RaggedRight\arraybackslash}X}
    Sekolah & Menetapkan kebijakan, jadwal, dukungan, dan tindak lanjut. \\
    Tentor & Menyusun materi, menganalisis kesalahan, dan menentukan remedial. \\
    Guru BK & Mendampingi motivasi dan pilihan program studi. \\
    Siswa & Menuntaskan materi, latihan, tryout, dan rencana perbaikan. \\
    StudentPro & Menyatukan aktivitas, hasil, umpan balik, dan monitoring. \\
  \end{tabularx}}
\end{frame}

% Slide 11 --------------------------------------------------------------------
\begin{frame}{Bab berikutnya menetapkan titik acuan}
  \kicker{Transisi ke Bab 2}
  \claim{Sebelum merancang latihan 2027, sekolah perlu memahami dengan tepat apa yang diukur UTBK 2026.}
  \vspace{0.45cm}
  \begin{itemize}
    \item Bab 2 menguraikan dua kelompok tes dan tujuh subtes.
    \item Jumlah soal, waktu, tujuan, serta tuntutan kognitif ditulis berdasarkan sumber resmi.
    \item Perubahan 2027 baru dinyatakan resmi setelah dipublikasikan oleh SNPMB.
  \end{itemize}
\end{frame}

% Slide 12 --------------------------------------------------------------------
\begin{frame}{Sumber resmi dan publik}
  \kicker{Rujukan Bab 1 - Bagian 1}
  {\color{StudentBlue}\bfseries\fontsize{16}{19}\selectfont 1. Panitia SNPMB}\par
  \vspace{0.08cm}
  {\fontsize{13}{16}\selectfont ``Framework UTBK-SNBT Tahun 2026.''}\par
  {\fontsize{12}{15}\selectfont\href{https://snpmb.id/fr/}{https://snpmb.id/fr/}}\par
  {\color{StudentGray}\fontsize{11.5}{14}\selectfont Sumber resmi struktur, tujuan tes, jumlah soal, waktu, dan tuntutan kompetensi.}\par
  \vspace{0.38cm}
  {\color{StudentBlue}\bfseries\fontsize{16}{19}\selectfont 2. DetikEdu - 26 Mei 2026}\par
  \vspace{0.08cm}
  {\fontsize{13}{16}\selectfont ``Skor LBI di Sertifikat UTBK 2026 Dipecah Saintek-Soshum, SNPMB Jelaskan Mengapa.''}\par
  {\fontsize{11.5}{14}\selectfont\href{https://www.detik.com/edu/seleksi-masuk-pt/d-8505258/}{detik.com/edu/seleksi-masuk-pt/d-8505258/}}\par
  {\color{StudentGray}\fontsize{11.5}{14}\selectfont Laporan konferensi pers yang mengutip pimpinan SNPMB 2026.}\par
\end{frame}

% Slide 13 --------------------------------------------------------------------
\begin{frame}{Sumber internal dan status TWK}
  \kicker{Rujukan Bab 1 - Bagian 2}
  {\color{StudentBlue}\bfseries\fontsize{16}{19}\selectfont 3. Dokumen internal StudentPro}\par
  \vspace{0.08cm}
  {\fontsize{13}{16}\selectfont ``LPJ Pembimbingan UTBK 2026.'' Penyusun: Balya Rochmadi.}\par
  {\color{StudentGray}\fontsize{11.5}{14}\selectfont Sumber rancangan program, target internal, dan evaluasi pembimbingan.}\par
  \vspace{0.55cm}
  {\color{StudentBlue}\bfseries\fontsize{16}{19}\selectfont 4. Rencana integrasi TWK ke LBI}\par
  \vspace{0.08cm}
  {\fontsize{13}{16}\selectfont Informasi perencanaan untuk mengantisipasi kemungkinan muatan wawasan kebangsaan dalam LBI.}\par
  {\color{StudentGray}\fontsize{11.5}{14}\selectfont Status: belum diposisikan sebagai ketentuan resmi sampai dokumen SNPMB 2027 diterbitkan.}\par
\end{frame}

\end{document}
